% ======================================================================
% 09_background.tex: CHAPTER 2, theoretical background
% WHAT TO WRITE
%   The theory the reader needs to follow your work: concepts, models, formulas, algorithms.
%   One \section per topic, one \subsection per sub-topic. Cite sources with \cite{key} (keys are in
%   backmatter/14_references.tex). Number important formulas with the equation environment and refer to them with \eqref{...}.
%   The last section shows how to cite and how to show code.
%
% HOW TO USE THIS FILE
%   - Replace every ... with your text. A new paragraph starts after one blank line.
%   - The title in \chapter{...} is printed as "CHAPTER n: TITLE" in capitals. Rename it freely.
%   - \section{...} and \subsection{...} are numbered automatically (1.1, 1.1.1). Rename, copy or delete them.
%   - \paragraph{Small heading:} gives a bold line of its own (the colon is required).
%   - The last section ("Template Examples") shows a figure, an equation and a table. Copy a block into your
%     text, edit it, then delete the "Template Examples" section when you no longer need it.
%   - Syntax for everything (lists, figures, tables, equations, citations): docs/LATEX_CHEATSHEET.md
%   - Labels (\label{...}) must be unique across the whole report.
% ======================================================================

\chapter{THEORETICAL BACKGROUND}
\label{chap:background}

\section{Topic A}
...

\subsection{Concept}
...

\subsection{Types and classification}
...

\subsection{Working principle}
...

\section{Topic B}
...

\subsection{Mathematical model}
...

\subsection{Notation}
...

\section{Topic C}
...

\subsection{Method}
...

\subsection{Algorithm}
...

\section{Citations and Code}
A reference is cited with a key from \texttt{backmatter/14\_references.tex}, for example \texttt{\textbackslash cite\{key\}} gives \cite{example_article}.
Source code is shown with the \texttt{pythoncode} environment:
\begin{pythoncode}
def f(x):
    return 2 * x
\end{pythoncode}

% ============================================================
% Template examples: copy these blocks into the relevant subsection.
% Replace placeholder text, captions, labels, and sample data as needed.
% ============================================================
\section{Template Examples}
\label{sec:template-ch2-background}

\subsection{Figure and Cross-Reference Example}
...

Figure~\ref{fig:template-ch2-background} illustrates ...
\begin{figure}[htbp]
    \centering
    % Replace the framed placeholder with your image, for example:
    % \includegraphics[width=0.75\linewidth]{figures/example.png}
    \fbox{%
        \begin{minipage}[c][4cm][c]{0.75\linewidth}
            \centering
            ...
        \end{minipage}%
    }
    \caption{...}
    \label{fig:template-ch2-background}
\end{figure}

\subsection{Equation Example}
...

An example relationship is given in Equation~\eqref{eq:template-ch2-background}:
\begin{equation}
    y(t) = a\,x(t) + b,
    \label{eq:template-ch2-background}
\end{equation}
where $x(t)$ is the input, $y(t)$ is the output, and $a$ and $b$ are
example coefficients. Replace this relationship and its definitions with ...

\subsection{Table Example}
...

Table~\ref{tab:template-ch2-background} presents ...
\begin{table}[htbp]
    \centering
    \caption{...}
    \label{tab:template-ch2-background}
    \begin{tabularx}{0.9\linewidth}{lXX}
        \toprule
        \textbf{Item} & \textbf{Description} & \textbf{Value / Unit} \\
        \midrule
        ... & ... & ... \\
        ... & ... & ... \\
        \bottomrule
    \end{tabularx}
\end{table}
